\documentclass[preview]{standalone}

\usepackage{amsmath}
\usepackage{amssymb}
\usepackage{stellar}
\usepackage{definitions}
\usepackage{bettelini}

\begin{document}

\id{probability-sums-random-variables}
\genpage

\section{Bernoulli sums}

\begin{snippettheorem}{binomial-distribution-as-bernoulli-sum}{Binomial distribution as a Bernoulli sum}
    Let \((\Omega,\powerset(\Omega),\mathbb P)\) be a
    \countableprobabilityspace, let \(n\in\naturalnumbers^\exceptzero\), let
    \(p\in[0,1]\), and let \(X_1,\ldots,X_n\) be
    \iiddiscreterandomvariables such that each \(X_i\) has
    \bernoullidistribution
    \[
        X_i\sim\operatorname{Bernoulli}(p)
    \]
    for every \(i\in\{1,\ldots,n\}\). Define
    \[
        S_n=\sum_{i=1}^n X_i.
    \]
    Then \(S_n\) has \binomialdistribution \(\operatorname{Binomial}(n,p)\).
\end{snippettheorem}

\begin{snippetproof}{binomial-distribution-as-bernoulli-sum-proof}{binomial-distribution-as-bernoulli-sum}{Binomial distribution as a Bernoulli sum}
    For \(k\in\{0,\ldots,n\}\), the \event \(\{S_n=k\}\) is the disjoint
    union of the events
    \[
        \{X_1=\varepsilon_1,\ldots,X_n=\varepsilon_n\},
        \qquad
        (\varepsilon_1,\ldots,\varepsilon_n)\in\{0,1\}^n,
        \qquad
        \sum_{i=1}^n\varepsilon_i=k.
    \]
    By independence, each such \event has probability
    \(p^k(1-p)^{n-k}\). The number of vectors
    \((\varepsilon_1,\ldots,\varepsilon_n)\in\{0,1\}^n\) with exactly \(k\)
    entries equal to \(1\) is \(\binom{n}{k}\). Hence
    \[
        \mathbb P(S_n=k)=\binom{n}{k}p^k(1-p)^{n-k}.
    \]
    This is the \randomvariabledensity of
    \(\operatorname{Binomial}(n,p)\).
\end{snippetproof}

\section{Sums and convolution}

\begin{snippetproposition}{sum-of-two-discrete-random-variables-density}{Density of a sum of two discrete random variables}
    Let \((\Omega,\powerset(\Omega),\mathbb P)\) be a
    \countableprobabilityspace. Let \(E,F\subseteq\realnumbers\) be finite or
    countably infinite \set[sets], and let
    \(X\colon\Omega\fromto E\) and \(Y\colon\Omega\fromto F\) be
    real-valued \discreterandomvariable[discrete random variables]. Define
    \[
        Z=X+Y,
        \qquad
        E+F=\{x+y\in\realnumbers\suchthat x\in E,\ y\in F\}.
    \]
    Then \(Z\colon\Omega\fromto E+F\) is a
    \discreterandomvariable and, for every \(z\in E+F\),
    \[
        p_Z(z)
        =
        \sum_{\substack{x\in E\\ z-x\in F}}p_{X,Y}(x,z-x)
        =
        \sum_{\substack{y\in F\\ z-y\in E}}p_{X,Y}(z-y,y).
    \]
\end{snippetproposition}

\begin{snippetproof}{sum-of-two-discrete-random-variables-density-proof}{sum-of-two-discrete-random-variables-density}{Density of a sum of two discrete random variables}
    Since \(E\) and \(F\) are finite or countably infinite, the \set
    \(E+F\) is finite or countably infinite. Thus \(Z\) is a
    \discreterandomvariable. For every \(z\in E+F\), the events
    \[
        \{X=x,\ Y=z-x\},
        \qquad x\in E,\ z-x\in F,
    \]
    are pairwise disjoint and their union is \(\{Z=z\}\). Therefore
    \[
        p_Z(z)=\mathbb P(Z=z)
        =
        \sum_{\substack{x\in E\\z-x\in F}}
        \mathbb P(X=x,\ Y=z-x)
        =
        \sum_{\substack{x\in E\\z-x\in F}}p_{X,Y}(x,z-x).
    \]
    The second formula is obtained by decomposing the same \event according
    to the value of \(Y\).
\end{snippetproof}

\begin{snippetcorollary}{discrete-sum-independent-convolution-formula}{Convolution formula for independent variables}
    Let \((\Omega,\powerset(\Omega),\mathbb P)\) be a
    \countableprobabilityspace. Let \(E,F\subseteq\realnumbers\) be finite or
    countably infinite \set[sets], and let
    \(X\colon\Omega\fromto E\) and \(Y\colon\Omega\fromto F\) be independent
    real-valued \discreterandomvariable[discrete random variables]. If
    \(Z=X+Y\), then, for every \(z\in E+F\),
    \[
        p_Z(z)
        =
        \sum_{\substack{x\in E\\ z-x\in F}}p_X(x)p_Y(z-x)
        =
        \sum_{\substack{y\in F\\ z-y\in E}}p_X(z-y)p_Y(y).
    \]
\end{snippetcorollary}

\begin{snippetproof}{discrete-sum-independent-convolution-formula-proof}{discrete-sum-independent-convolution-formula}{Convolution formula for independent variables}
    By independence, the \jointdensity satisfies
    \[
        p_{X,Y}(x,y)=p_X(x)p_Y(y)
    \]
    for every \((x,y)\in E\cartesianprod F\). Substituting this identity in
    the density formula for \(X+Y\) gives the result.
\end{snippetproof}

\begin{snippet}{discrete-convolution-interpretation}
    The preceding formula says that the \randomvariabledensity of the sum of
    two independent real-valued \discreterandomvariable[discrete random
    variables] is the \discreteconvolution of their densities.
\end{snippet}

\begin{snippetcorollary}{sum-independent-poisson-random-variables}{Sum of independent Poisson random variables}
    Let \((\Omega,\powerset(\Omega),\mathbb P)\) be a
    \countableprobabilityspace. Let \(X,Y\colon\Omega\fromto\naturalnumbers\)
    be independent \discreterandomvariable[discrete random variables]. If
    \(X\) has \poissondistribution \(\operatorname{Poisson}(\lambda)\) and
    \(Y\) has \poissondistribution \(\operatorname{Poisson}(\mu)\), with
    \(\lambda,\mu>0\), then \(X+Y\) has \poissondistribution
    \(\operatorname{Poisson}(\lambda+\mu)\).
\end{snippetcorollary}

\begin{snippetproof}{sum-independent-poisson-random-variables-proof}{sum-independent-poisson-random-variables}{Sum of independent Poisson random variables}
    Let \(Z=X+Y\). For every \(z\in\naturalnumbers\), the
    \convolutionformula gives
    \begin{align*}
        p_Z(z)
        &=
        \sum_{x=0}^z
        e^{-\lambda}\frac{\lambda^x}{x!}
        e^{-\mu}\frac{\mu^{z-x}}{(z-x)!}
        \\
        &=
        e^{-(\lambda+\mu)}
        \sum_{x=0}^z
        \frac{\lambda^x\mu^{z-x}}{x!(z-x)!}
        \\
        &=
        e^{-(\lambda+\mu)}
        \frac{1}{z!}
        \sum_{x=0}^z
        \binom{z}{x}\lambda^x\mu^{z-x}
        \\
        &=
        e^{-(\lambda+\mu)}
        \frac{(\lambda+\mu)^z}{z!}.
    \end{align*}
    This is the \randomvariabledensity of
    \(\operatorname{Poisson}(\lambda+\mu)\).
\end{snippetproof}

\end{document}
